% ============================================================
% Section 125: 价带顶空穴性质计算
% ============================================================
\section{半导体物理：价带顶空穴性质计算}
\label{sec:hole-properties}

\begin{modelbox}[题目：价带顶附近空穴的准粒子性质]
在价带顶附近波矢为 $k$ 的一个态中，电子能量由下式给出：
\[
E=-10^{-37}k^2\quad(\text{J}).
\]
若把一个电子从 $k_1=1\times10^9\,\text{m}^{-1}$ 态上除去，而所有其他态均被电子占据。试计算：

(1) 空穴的有效质量；

(2) 空穴的波矢；

(3) 空穴的速度；

(4) 空穴的能量（以价带顶为参考点），并在 $E_h$--$k_h$ 能带图上示出该空穴。
\end{modelbox}

\begin{tipbox}[考点快速分析]
\begin{itemize}
    \item \textbf{空穴是准粒子}：满带中缺失电子的等价描述，带正电、正有效质量
    \item \textbf{对应关系}：$m_h^*=-m_e^*$，$k_h=-k_e$，$v_h=v_e$，$E_h=-E_e$（以价带顶为零点）
    \item \textbf{电子有效质量}：由色散关系 $E(k)$ 的二阶导数确定
    \item \textbf{群速度}：$v=\hbar^{-1}\dd{E}{k}$
\end{itemize}
\end{tipbox}

\begin{derivationbox}[标准解题过程]
\textbf{(1) 空穴的有效质量 $m_h^*$}

电子在价带顶附近的能量色散关系为
\[
E(k)=\frac{\hbar^2k^2}{2m_e^*}=-10^{-37}k^2\quad(\text{J}).
\]
由此可得电子的有效质量：
\[
m_e^*=\frac{\hbar^2}{2\times(-10^{-37})}=-\frac{(1.055\times10^{-34})^2}{2\times10^{-37}}\approx-5.57\times10^{-32}\,\text{kg}\approx-0.061\,m_0.
\]
空穴的有效质量是缺失电子有效质量的相反数：
\begin{equation}
\boxed{m_h^*=-m_e^*\approx0.061\,m_0.}
\end{equation}

\textbf{(2) 空穴的波矢 $k_h$}

空穴的波矢等于缺失电子波矢的相反数：
\begin{equation}
\boxed{k_h=-k_1=-1\times10^9\,\text{m}^{-1}.}
\end{equation}

\textbf{(3) 空穴的速度 $v_h$}

电子在 $k_1$ 处的群速度：
\[
v_e=\frac{1}{\hbar}\dd{E}{k}\bigg|_{k=k_1}=\frac{1}{\hbar}\cdot2\times(-10^{-37})\cdot k_1=-\frac{2\times10^{-37}\times10^9}{1.055\times10^{-34}}\approx-1.90\times10^6\,\text{m/s}.
\]
空穴的速度等于该处电子的速度：
\begin{equation}
\boxed{v_h=v_e\approx-1.90\times10^6\,\text{m/s}.}
\end{equation}

\textbf{(4) 空穴的能量 $E_h$}

以价带顶（$k=0$）为能量参考点，$E(0)=0$。在 $k_1$ 处电子的能量为
\[
E(k_1)=-10^{-37}\times(10^9)^2=-1\times10^{-19}\,\text{J}.
\]
空穴的能量定义为缺失电子能量的相反数（以价带顶为零点，向下为正）：
\[
E_h=-E(k_1)=1\times10^{-19}\,\text{J}=\frac{1\times10^{-19}}{1.602\times10^{-19}}\approx0.624\,\text{eV}.
\]
\begin{equation}
\boxed{E_h\approx0.62\,\text{eV}.}
\end{equation}

在空穴能带图（$E_h$--$k_h$，纵轴为空穴能量，向下为正）上：价带顶在 $k_h=0$、$E_h=0$；该空穴位于 $k_h=-1\times10^9\,\text{m}^{-1}$、$E_h\approx0.62\,\text{eV}$ 处。
\end{derivationbox}

\begin{formulabox}[答案汇总]
\begin{center}
\begin{tabular}{ll}
\toprule
\textbf{物理量} & \textbf{数值} \\
\midrule
(1) 空穴有效质量 $m_h^*$ & $0.061\,m_0$ \\
(2) 空穴波矢 $k_h$ & $-1\times10^9\,\text{m}^{-1}$ \\
(3) 空穴速度 $v_h$ & $-1.90\times10^6\,\text{m/s}$ \\
(4) 空穴能量 $E_h$ & $0.62\,\text{eV}$ \\
\bottomrule
\end{tabular}
\end{center}
\end{formulabox}

\begin{insightbox}[物理图像]
\begin{itemize}
    \item 空穴是满带中缺失电子形成的准粒子，其有效质量、波矢、速度和能量均与缺失的电子有确定的关系。
    \item 空穴有效质量为正值，便于用类自由粒子模型描述近满带的导电行为。
    \item 空穴能量惯例向下增加，空穴倾向于占据能量较低的状态（即电子带顶附近）。这符合``空穴浮向价带顶''的直观图像。
\end{itemize}
\end{insightbox}
